% claim-ref: mosquitofish#livebearing-sexing
% claim-ref: mosquitofish#stored-sperm
% claim-ref: mosquitofish#variable-timing
% claim-ref: mosquitofish#responsible-sequence
% claim-ref: mosquitofish#capacity-stop
% claim-ref: mosquitofish#no-release
\CompanionHeader{REPRODUCTION}{Livebearing and capacity}{Western mosquitofish / responsible observation}{Fry are born free-swimming, not as an externally laid egg clutch. [MF-MO]}
\Context{Breeding is optional. Inclusion in this binder does not authorize acquisition, transfer or release, or adding fish to a shrimp nursery or Kuhli community.}
\Step{1}{Decide and arrange surplus care}{Confirm spare prepared capacity and a district-coordinated destination before encouraging reproduction. Never use starvation, poor water or deliberate predation to regulate numbers. [MF-VET; proposed]}
\Step{2}{Verify readiness and adult sex}{Healthy adults need a proven suitable habitat. Mature males have an elongated modified anal fin (gonopodium); females generally have a fan-shaped anal fin and larger body. Juvenile size or a dark spot alone is not reliable. [MF-MO]}
\Step{3}{Allow observation, not forced production}{Females can store sperm; separation from males may not stop subsequent births. Observe courtship and births without forced pairing or routine confinement in a tiny breeder box. [MF-FWS; proposed]}
\Step{4}{Keep timing in scope}{A federal synthesis reports 21--28 days gestation; brood size and season vary. A laboratory abstract reports mean first reproduction at 191 days at 20 C versus 56 days at 30 C. These are experimental endpoints, NOT heater targets or monthly-brood promises. Do not heat to force reproduction. [MF-FWS; MF-VON]}
\Step{5}{Support fry safely}{Adults may eat young. Provide fine refuge, safe intakes and accessible small food while preserving open water. A separate nursery needs established biofiltration and suitable matched water before use; cover alone cannot guarantee survival. [MF-SM; MF-VET; proposed]}
\Step{6}{Record milestones and stop cues}{Record births, feeding, growth and later reproductive adults separately. Review adults/fry and capacity weekly during breeding. Harassment, poor feeding, unusual losses or unstable water pause expansion/transfer while conditions and health are reviewed. [MF-VET; proposed]}
\Step{7}{Contain and coordinate}{Never release to natural waters or drains. Ask San Mateo district about unwanted-fish collection: 650-344-8592. Confirm arrangements before breeding and verify current local permissions before any placement or transfer. Avoid amphibian-conservation ponds; predation and ecological harm are documented. [MF-RETURN; MF-FWS]}
\Note{Observation record}{Births / fry feeding / growth / adults / water: \hrulefill\par Capacity / surplus destination / next review: \hrulefill}
\Sources{Clickable evidence: \href{https://mdc.mo.gov/discover-nature/field-guide/western-mosquitofish}{MF-MO}; \href{https://www.fws.gov/sites/default/files/documents/Ecological-Risk-Screening-Summary-Western-Mosquitofish.pdf}{MF-FWS}; \href{https://digitalcommons.usu.edu/wats_facpub/75/}{MF-VON}; \href{https://www.smcmvcd.org/mosquitofish}{MF-SM}; \href{https://www.msdvetmanual.com/exotic-and-laboratory-animals/aquarium-fish/management-of-aquarium-fish}{MF-VET}; \href{https://www.smcmvcd.org/request-mosquitofish-for-your-backyard-pond-or-water-garden}{MF-RETURN}. Source passages, dates and applicability: adjacent YAML. Planning only; actual habitat unverified.}{mosquitofish / reproduction}
